\begin{thebibliography}{99}
\footnotesize\itemsep=0pt

\bibitem{Ar1}
Aratyn H., Gomes J.F., Lobo G.V., Zimerman A.H., Why is my rational
 Painlev\'e~V solution not unique?,
 \href{http://arxiv.org/abs/2307.07825}{arXiv:2307.07825}.

\bibitem{Ar3}
Aratyn H., Gomes J.F., Lobo G.V., Zimerman A.H., Extended symmetry of higher
 Painlev\'e equations of even periodicity and their rational solutions,
 \href{https://doi.org/10.3390/math12233701}{\textit{Mathematics}} \textbf{12}
 (2024), 3701, 26~pages,
 \href{http://arxiv.org/abs/2409.03534}{arXiv:2409.03534}.

\bibitem{Ar2}
Aratyn H., Gomes J.F., Lobo G.V., Zimerman A.H., Two-fold degeneracy of a class
 of rational {P}ainlev\'e~{V} solutions,
 \href{https://doi.org/10.46298/ocnmp.13201}{\textit{Open Commun. Nonlinear
 Math. Phys.}} (2024), spec. issue 2, 28--45,
 \href{http://arxiv.org/abs/2310.01585}{arXiv:2310.01585}.

\bibitem{BalBerBo}
Balogh F., Bertola M., Bothner T., Hankel determinant approach to generalized
 {V}orob'ev--{Y}ablonski polynomials and their roots,
 \href{https://doi.org/10.1007/s00365-016-9328-4}{\textit{Constr. Approx.}}
 \textbf{44} (2016), 417--453,
 \href{http://arxiv.org/abs/1504.00440}{arXiv:1504.00440}.

\bibitem{Deano}
Barhoumi A., Celsus A.F., Dea\~{n}o A., Global-phase portrait and large-degree
 asymptotics for the {K}issing polynomials,
 \href{https://doi.org/10.1111/sapm.12387}{\textit{Stud. Appl. Math.}}
 \textbf{147} (2021), 448--526,
 \href{http://arxiv.org/abs/2008.08724}{arXiv:2008.08724}.

\bibitem{BertolaBothner}
Bertola M., Bothner T., Zeros of large degree {V}orob'ev--{Y}ablonski
 polynomials via a {H}ankel determinant identity,
 \href{https://doi.org/10.1093/imrn/rnu239}{\textit{Int. Math. Res. Not.}}
 \textbf{2015} (2015), 9330--9399,
 \href{http://arxiv.org/abs/1401.1408}{arXiv:1401.1408}.

\bibitem{Bertola:Semiiso}
Bertola M., Eynard B., Harnad J., Semiclassical orthogonal polynomials, matrix
 models and isomonodromic tau functions,
 \href{https://doi.org/10.1007/s00220-005-1505-4}{\textit{Comm. Math. Phys.}}
 \textbf{263} (2006), 401--437,
 \href{http://arxiv.org/abs/nlin.SI/0410043}{arXiv:nlin.SI/0410043}.

\bibitem{BertoKatTov}
Bertola M., Katsevich A., Tovbis A., Singular value decomposition of a finite
 {H}ilbert transform defined on several intervals and the interior problem of
 tomography: the {R}iemann--{H}ilbert problem approach,
 \href{https://doi.org/10.1002/cpa.21547}{\textit{Comm. Pure Appl. Math.}}
 \textbf{69} (2016), 407--477,
 \href{http://arxiv.org/abs/1402.0216}{arXiv:1402.0216}.

\bibitem{Bertola:wj}
Bertola M., Lee S.Y., First colonization of a hard-edge in random matrix
 theory, \href{https://doi.org/10.1007/s00365-009-9052-4}{\textit{Constr.
 Approx.}} \textbf{31} (2010), 231--257,
 \href{http://arxiv.org/abs/0804.1111}{arXiv:0804.1111}.

\bibitem{BertolaAdvances}
Bertola M., Mo M.Y., Commuting difference operators, spinor bundles and the
 asymptotics of orthogonal polynomials with respect to varying complex
 weights, \href{https://doi.org/10.1016/j.aim.2008.09.001}{\textit{Adv.
 Math.}} \textbf{220} (2009), 154--218,
 \href{http://arxiv.org/abs/math-ph/0605043}{arXiv:math-ph/0605043}.

\bibitem{BleherIts}
Bleher P., Its A., Double scaling limit in the random matrix model: the
 {R}iemann--{H}ilbert approach,
 \href{https://doi.org/10.1002/cpa.10065}{\textit{Comm. Pure Appl. Math.}}
 \textbf{56} (2003), 433--516,
 \href{http://arxiv.org/abs/math-ph/0201003}{arXiv:math-ph/0201003}.

\bibitem{Borodin}
Borodin A., Forrester P.J., Increasing subsequences and the hard-to-soft edge
 transition in matrix ensembles,
 \href{https://doi.org/10.1088/0305-4470/36/12/307}{\textit{J.~Phys.~A}}
 \textbf{36} (2003), 2963--2981,
 \href{http://arxiv.org/abs/math-ph/0205007}{arXiv:math-ph/0205007}.

\bibitem{Ref1_2}
Bothner T., Miller P.D., Rational solutions of the {P}ainlev\'e-{III} equation:
 large parameter asymptotics,
 \href{https://doi.org/10.1007/s00365-019-09463-4}{\textit{Constr. Approx.}}
 \textbf{51} (2020), 123--224,
 \href{http://arxiv.org/abs/1808.01421}{arXiv:1808.01421}.

\bibitem{Ref1_1}
Bothner T., Miller P.D., Sheng Y., Rational solutions of the {P}ainlev\'e-{III}
 equation, \href{https://doi.org/10.1111/sapm.12220}{\textit{Stud. Appl.
 Math.}} \textbf{141} (2018), 626--679,
 \href{http://arxiv.org/abs/1801.04360}{arXiv:1801.04360}.

\bibitem{Ref1_3}
Buckingham R.J., Large-degree asymptotics of rational {P}ainlev\'e-{IV}
 functions associated to generalized {H}ermite polynomials,
 \href{https://doi.org/10.1093/imrn/rny172}{\textit{Int. Math. Res. Not.}}
 \textbf{2020} (2020), 5534--5577,
 \href{http://arxiv.org/abs/1706.09005}{arXiv:1706.09005}.

\bibitem{BuckMill12}
Buckingham R.J., Miller P.D., The sine-{G}ordon equation in the semiclassical
 limit: critical behavior near a separatrix,
 \href{https://doi.org/10.1007/s11854-012-0041-3}{\textit{J.~Anal. Math.}}
 \textbf{118} (2012), 397--492,
 \href{http://arxiv.org/abs/1106.5716}{arXiv:1106.5716}.

\bibitem{BuckMill14}
Buckingham R.J., Miller P.D., Large-degree asymptotics of rational
 {P}ainlev\'e-{II} functions: noncritical behaviour,
 \href{https://doi.org/10.1088/0951-7715/27/10/2489}{\textit{Nonlinearity}}
 \textbf{27} (2014), 2489--2578,
 \href{http://arxiv.org/abs/1310.2276}{arXiv:1310.2276}.

\bibitem{BuckMill15}
Buckingham R.J., Miller P.D., Large-degree asymptotics of rational
 {P}ainlev\'e-{II} functions: critical behaviour,
 \href{https://doi.org/10.1088/0951-7715/28/6/1539}{\textit{Nonlinearity}}
 \textbf{28} (2015), 1539--1596,
 \href{http://arxiv.org/abs/1406.0826}{arXiv:1406.0826}.

\bibitem{BuckMill22}
Buckingham R.J., Miller P.D., Large-degree asymptotics of rational
 {P}ainlev\'e-{IV} solutions by the isomonodromy method,
 \href{https://doi.org/10.1007/s00365-022-09586-1}{\textit{Constr. Approx.}}
 \textbf{56} (2022), 233--443,
 \href{http://arxiv.org/abs/2008.00600}{arXiv:2008.00600}.

\bibitem{Clakxon_IV}
Clarkson P.A., The fourth {P}ainlev\'e equation and associated special
 polynomials, \href{https://doi.org/10.1063/1.1603958}{\textit{J.~Math.
 Phys.}} \textbf{44} (2003), 5350--5374.

\bibitem{ClaxonIII}
Clarkson P.A., The third {P}ainlev\'e equation and associated special
 polynomials,
 \href{https://doi.org/10.1088/0305-4470/36/36/306}{\textit{J.~Phys.~A}}
 \textbf{36} (2003), 9507--9532.

\bibitem{ClarksonV}
Clarkson P.A., Special polynomials associated with rational solutions of the
 {P}ainlev\'e equations and applications to soliton equations,
 \href{https://doi.org/10.1007/BF03321618}{\textit{Comput. Methods Funct.
 Theory}} \textbf{6} (2006), 329--401,
 \href{http://arxiv.org/abs/nlin.SI/0605034}{arXiv:nlin.SI/0605034}.

\bibitem{ClarksonDunning}
Clarkson P.A., Dunning C., Rational solutions of the fifth {P}ainlev\'e
 equation. {G}eneralized {L}aguerre polynomials,
 \href{https://doi.org/10.1111/sapm.12649}{\textit{Stud. Appl. Math.}}
 \textbf{152} (2024), 453--507,
 \href{http://arxiv.org/abs/2304:01579}{arXiv:2304:01579}.

\bibitem{Clarxon}
Clarkson P.A., Mansfield E.L., The second {P}ainlev\'e equation, its hierarchy
 and associated special polynomials,
 \href{https://doi.org/10.1088/0951-7715/16/3/201}{\textit{Nonlinearity}}
 \textbf{16} (2003), R1--R26.

\bibitem{Deiftbook}
Deift P., Orthogonal polynomials and random matrices: a {R}iemann--{H}ilbert
 approach, \textit{Courant Lect. Notes Math.}, Vol.~3, New York University,
 Courant Institute of Mathematical Sciences, New York, 1999.

\bibitem{DKMVZ}
Deift P., Kriecherbauer T., McLaughlin K.T.-R., Venakides S., Zhou X., Uniform
 asymptotics for polynomials orthogonal with respect to varying exponential
 weights and applications to universality questions in random matrix theory,
 \href{https://doi.org/10.1002/(SICI)1097-0312(199911)52:11<1335::AID-CPA1>3.0.CO;2-1}{\textit{Comm.
 Pure Appl. Math.}} \textbf{52} (1999), 1335--1425.

\bibitem{DZ}
Deift P., Zhou X., A steepest descent method for oscillatory
 {R}iemann--{H}ilbert problems. {A}symptotics for the {MK}d{V} equation,
 \href{https://doi.org/10.2307/2946540}{\textit{Ann. of Math.}} \textbf{137}
 (1993), 295--368.

\bibitem{FIK}
Fokas A.S., Its A.R., Kitaev A.V., Discrete {P}ainlev\'e equations and their
 appearance in quantum gravity,
 \href{https://doi.org/10.1007/BF02102066}{\textit{Comm. Math. Phys.}}
 \textbf{142} (1991), 313--344.

\bibitem{Gakhov}
Gakhov F.D., Boundary value problems, Dover Publications, Inc., New York, 1990.

\bibitem{Gambier}
Gambier B., Sur les \'equations diff\'erentielles du second ordre et du premier
 degr\'e dont l'int\'egrale g\'en\'erale est a points critiques fixes,
 \href{https://doi.org/10.1007/BF02393211}{\textit{Acta Math.}} \textbf{33}
 (1910), 1--55.

\bibitem{HasMc}
Hastings S.P., McLeod J.B., A boundary value problem associated with the second
 {P}ainlev\'e transcendent and the {K}orteweg-de\thinspace {V}ries equation,
 \href{https://doi.org/10.1007/BF00283254}{\textit{Arch. Rational Mech.
 Anal.}} \textbf{73} (1980), 31--51.

\bibitem{JMU1}
Jimbo M., Miwa T., Ueno K., Monodromy preserving deformation of linear ordinary
 differential equations with rational coefficients.~{I}. {G}eneral theory and
 {$\tau $}-function,
 \href{https://doi.org/10.1016/0167-2789(81)90013-0}{\textit{Phys.~D}}
 \textbf{2} (1981), 306--352.

\bibitem{JMU2}
Jimbo M., Miwa T., Monodromy preserving deformation of linear ordinary
 differential equations with rational coefficients.~{II},
 \href{https://doi.org/10.1016/0167-2789(81)90021-X}{\textit{Phys.~D}}
 \textbf{2} (1981), 407--448.

\bibitem{JoshiKajiwara}
Joshi N., Kajiwara K., Mazzocco M., Generating function associated with the
 determinant formula for the solutions of the {P}ainlev\'e~{II} equation,
 \textit{Ast\'erisque} \textbf{297} (2004), 67--78,
 \href{http://arxiv.org/abs/nlin.SI/0406035}{arXiv:nlin.SI/0406035}.

\bibitem{KajiwaraOhta}
Kajiwara K., Ohta Y., Determinant structure of the rational solutions for the
 {P}ainlev\'e~{II} equation,
 \href{https://doi.org/10.1063/1.531648}{\textit{J.~Math. Phys.}} \textbf{37}
 (1996), 4693--4704,
 \href{http://arxiv.org/abs/solv-int/9607002}{arXiv:solv-int/9607002}.

\bibitem{Kitaev94}
Kitaev A.V., Law C.K., McLeod J.B., Rational solutions of the fifth
 {P}ainlev\'e equation, \textit{Differential Integral Equations} \textbf{7}
 (1994), 967--1000.

\bibitem{Marcellan0}
Marcell\'an F., Rocha I.A., On semiclassical linear functionals: integral
 representations,
 \href{https://doi.org/10.1016/0377-0427(93)E0248-K}{\textit{J.~Comput. Appl.
 Math.}} \textbf{57} (1995), 239--249.

\bibitem{Marcellan}
Marcell\'an F., Rocha I.A., Complex path integral representation for
 semiclassical linear functionals,
 \href{https://doi.org/10.1006/jath.1998.3190}{\textit{J.~Approx. Theory}}
 \textbf{94} (1998), 107--127.

\bibitem{Maroni}
Maroni P., Prol\'egom\`enes \`a l'\'etude des polyn\^{o}mes orthogonaux
 semi-classiques, \href{https://doi.org/10.1007/BF01773932}{\textit{Ann. Mat.
 Pura Appl.}} \textbf{149} (1987), 165--184.

\bibitem{MasoeroIV}
Masoero D., Roffelsen P., Poles of {P}ainlev\'e~{IV} rationals and their
 distribution, \href{https://doi.org/10.3842/SIGMA.2018.002}{\textit{SIGMA}}
 \textbf{14} (2018), 002, 49~pages,
 \href{http://arxiv.org/abs/1707.05222}{arXiv:1707.05222}.

\bibitem{MasoeroIV2}
Masoero D., Roffelsen P., Roots of generalised {H}ermite polynomials when both
 parameters are large,
 \href{https://doi.org/10.1088/1361-6544/abdd93}{\textit{Nonlinearity}}
 \textbf{34} (2021), 1663--1732,
 \href{http://arxiv.org/abs/1907.08552}{arXiv:1907.08552}.

\bibitem{Masuda}
Masuda T., Some properties of zeros of the umemura polynomials of the fifth
 {P}ainlev\'e equation,
 \href{https://doi.org/10.1111/sapm.70109}{\textit{Stud. Appl. Math.}}
 \textbf{155} (2025), e70109, 34~pages.

\bibitem{MR1942391}
Masuda T., Ohta Y., Kajiwara K., A determinant formula for a class of rational
 solutions of {P}ainlev\'e~{V} equation,
 \href{https://doi.org/10.1017/S0027763000008321}{\textit{Nagoya Math.~J.}}
 \textbf{168} (2002), 1--25,
 \href{http://arxiv.org/abs/nlin.SI/0101056}{arXiv:nlin.SI/0101056}.

\bibitem{Marta}
Mazzocco M., Rational solutions of the {P}ainlev\'e~{VI} equation,
 \href{https://doi.org/10.1088/0305-4470/34/11/320}{\textit{J.~Phys.~A}}
 \textbf{34} (2001), 2281--2294,
 \href{http://arxiv.org/abs/nlin.SI/0007036}{arXiv:nlin.SI/0007036}.

\bibitem{McCoyTracyWu1}
McCoy B.M., Tracy C.A., Wu T.T., Connection between the {K}d{V} equation and
 the two-dimensional {I}sing model,
 \href{https://doi.org/10.1016/0375-9601(77)90613-2}{\textit{Phys. Lett.~A}}
 \textbf{61} (1977), 283--284.

\bibitem{McCoyTracyWu2}
McCoy B.M., Tracy C.A., Wu T.T., Painlev\'e functions of the third kind,
 \href{https://doi.org/10.1063/1.523367}{\textit{J.~Math. Phys.}} \textbf{18}
 (1977), 1058--1092.

\bibitem{Ref1_4}
Miller P.D., Sheng Y., Rational solutions of the {P}ainlev\'e-{II} equation
 revisited, \href{https://doi.org/10.3842/SIGMA.2017.065}{\textit{SIGMA}}
 \textbf{13} (2017), 065, 29~pages,
 \href{http://arxiv.org/abs/1704.04851}{arXiv:1704.04851}.

\bibitem{MR1650509}
Noumi M., Yamada Y., Umemura polynomials for the {P}ainlev\'e~{V} equation,
 \href{https://doi.org/10.1016/S0375-9601(98)00625-2}{\textit{Phys. Lett.~A}}
 \textbf{247} (1998), 65--69.

\bibitem{Novok14}
Novokshenov V.Yu., Shchelkonogov A.A., Double scaling limit in the
 {P}ainlev\'e~{IV} equation and asymptotics of the {O}kamoto polynomials, in
 Spectral {T}heory and {D}ifferential {E}quations, \textit{Amer. Math. Soc.
 Transl. Ser.~2}, Vol. 233,
 \href{https://doi.org/10.1090/trans2/233/12}{American Mathematical Society},
 Providence, RI, 2014, 199--210.

\bibitem{Okamoto1}
Okamoto K., Studies on the {P}ainlev\'e equations.~{III}. {S}econd and fourth
 {P}ainlev\'e equations,~{$P_{{\rm II}}$} and~{$P_{{\rm IV}}$},
 \href{https://doi.org/10.1007/BF01458459}{\textit{Math. Ann.}} \textbf{275}
 (1986), 221--255.

\bibitem{Okamoto2}
Okamoto K., Studies on the {P}ainlev\'e equations.~{II}. {F}ifth {P}ainlev\'e
 equation~{$P_{\rm V}$},
 \href{https://doi.org/10.4099/math1924.13.47}{\textit{Japan.~J. Math.
 (N.S.)}} \textbf{13} (1987), 47--76.

\bibitem{Okamoto3}
Okamoto K., Studies on the {P}ainlev\'e equations.~{IV}. {T}hird {P}ainlev\'e
 equation~{$P_{{\rm III}}$}, \textit{Funkcial. Ekvac.} \textbf{30} (1987),
 305--332.

\bibitem{DLMF}
Olver F.W.J., Olde~Daalhuis A.B., Lozier D.W., Schneider B.I., Boisvert R.F.,
 Clark C.W., Miller B.R., Saunders B.V., Cohl H.S., McClain M.A. (Editors),
 NIST Digital Library of Mathematical Functions, 2025, Release 1.2.4.

\bibitem{Painleve}
Painlev\'e P., Sur les \'equations diff\'erentielles du second ordre et d'ordre
 sup\'erieur dont l'int\'egrale g\'en\'erale est uniforme,
 \href{https://doi.org/10.1007/BF02419020}{\textit{Acta Math.}} \textbf{25}
 (1902), 1--85.

\bibitem{Shohat}
Shohat J., A differential equation for orthogonal polynomials,
 \href{https://doi.org/10.1215/S0012-7094-39-00534-X}{\textit{Duke Math.~J.}}
 \textbf{5} (1939), 401--417.

\bibitem{Strebel}
Strebel K., Quadratic differentials, \textit{Ergeb. Math. Grenzgeb.~(3)},
 Vol.~5, \href{https://doi.org/10.1007/978-3-662-02414-0}{Springer}, Berlin,
 1984.

\bibitem{TracyWidom94}
Tracy C.A., Widom H., Level-spacing distributions and the {A}iry kernel,
 \href{https://doi.org/10.1007/BF02100489}{\textit{Comm. Math. Phys.}}
 \textbf{159} (1994), 151--174,
 \href{http://arxiv.org/abs/hep-th/9211141}{arXiv:hep-th/9211141}.

\bibitem{Umemura1}
Umemura H., Special polynomials associated with the {P}ainlev\'e equations~{I},
 \href{https://doi.org/10.5802/afst.1657}{\textit{Ann. Fac. Sci. Toulouse
 Math.~(6)}} \textbf{29} (2020), 1063--1089.

\bibitem{VanlessenStrong}
Vanlessen M., Strong asymptotics of {L}aguerre-type orthogonal polynomials and
 applications in random matrix theory,
 \href{https://doi.org/10.1007/s00365-005-0611-z}{\textit{Constr. Approx.}}
 \textbf{25} (2007), 125--175,
 \href{http://arxiv.org/abs/math.CA/0504604}{arXiv:math.CA/0504604}.

\bibitem{Vorobev}
Vorob'ev A.P., On the rational solutions of the second {P}ainlev\'e equation,
 \textit{Differ. Equ.} \textbf{1} (1965), 79--81.

\bibitem{Yablonskii}
Yablonskii A.I., Residues of poles of solutions of Painlev\'e's second
 equation, \textit{Dokl. Akad. Nauk BSSR} \textbf{4} (1960), 47--50.

\end{thebibliography}
